% !TeX program = xelatex
\documentclass[11pt,a4paper]{article}
\usepackage[margin=23mm]{geometry}
\usepackage{fontspec}
\setmainfont{Latin Modern Roman}
\usepackage{amsmath,amssymb,mathtools}
\usepackage{xcolor,enumitem,booktabs,longtable,array,xurl,hyperref}
\definecolor{muted}{HTML}{53636C}
\hypersetup{colorlinks=true,urlcolor=blue,linkcolor=blue}
\setlength{\parindent}{0pt}
\setlength{\parskip}{5pt}
\newcommand{\R}{\mathbb R}
\newcommand{\E}{\mathbb E}
\newcommand{\N}{\mathbb N}
\newcommand{\ind}{\mathbf 1}
\newcommand{\Var}{\operatorname{Var}}
\newcommand{\Cov}{\operatorname{Cov}}
\newcommand{\supp}{\operatorname{supp}}
\DeclareMathOperator*{\argmin}{arg\,min}
\DeclareMathOperator*{\argmax}{arg\,max}
\newcommand{\entrymeta}[3]{{\sffamily\small\color{muted}\textbf{\texttt{#1}}\quad|\quad #2\par #3\par}}
\newcommand{\Problemref}[1]{\texttt{#1}}
\newcommand{\JELref}[2]{\texttt{#2} (JEL \texttt{#1})}
\begin{document}
\section*{Core-market fragility}
\textbf{Problem ID:} \texttt{G18-100}\quad \textbf{JEL:} \texttt{G18}\par
% BEGIN ECONOMICS STATEMENT
\label{problem:OP-0317}
{\sffamily\small\color{muted}Original subject group: Asset pricing and financial markets\par}
\label{definitions:problem:30007599}
\textbf{Audit classification:} Published research agenda. The official agenda and BIS open-question section support market-resilience work. The robust architecture loss function is editorial; known compact optimization existence is not an unresolved theorem.
\subsection*{Definitions and precise research target}
Fix a government-bond market with dealers, leveraged funds, customer orders, collateral and a specified settlement network. A feasible architecture \(a\in A\) specifies trading access, central clearing, netting periods, collateral eligibility and margin schedules. At each date agents choose portfolios and liquidity buffers optimally, taking these rules into account. Let \(L(a,\omega)\) be realized social loss in shock state \(\omega\), combining execution costs, failed settlement, default externalities and real-economy disruption, with its components separately measured. Let \(C(a)\) be the ordinary resource and compliance cost. For a publicly defined uncertainty set \(\mathcal P\) of shock distributions, seek
\[J(a)=C(a)+\sup_{P\in\mathcal P}E_P[L(a,\omega)],\qquad J_*=\inf_{a\in A}J(a).\]
The research task is to construct an empirically disciplined loss function and characterize designs that remain resilient when normally liquid assets experience simultaneous selling and funding shocks. Specify institutional budget, legal and incentive constraints defining \(A\). Account for changes in leverage and liquidity provision induced by the architecture, rather than holding agent behavior fixed. Report the frontier between routine intermediation costs and tail losses, and the robustness of optimal designs to uncertain market-impact elasticities and shock dependence. A design optimal under one calibrated stress episode need not solve the wider market-resilience question.

\textbf{Model-specific implementation requirements (editorial).} This is an agenda for constructing and comparing empirically justified models, not a single calibrated optimization instance. Each implementation must declare the bond market, participant and affected-population boundaries, stress horizon and baseline settlement network. Specify whether architectures are static rules or history-dependent policies, their domains and information restrictions, and how each architecture changes obligations, netting sets and default arrangements. Declare a topology on \(A\) when invoking compactness or continuity; neither is imposed on the general agenda.

Define the dynamic equilibrium correspondence from preferences, endowments, financing constraints, information, beliefs, policy timing, price formation and market clearing. Establish or explicitly assume nonemptiness for the comparisons made, and declare an equilibrium selection rule or an explicit quantification over permitted equilibria. Individual portfolio optimality alone does not define \(L\). The displayed notation suppresses fixed model primitives and an explicitly declared equilibrium convention. State whether \(P\) indexes only exogenous innovations, whether agents know that law or face ambiguity, and how endogenous sales and funding outcomes are generated. If behavioral beliefs, equilibrium selection or loss depend on \(P\), retain that dependence explicitly in the loss inside the expectation; if shock laws depend on \(a\), specify the corresponding uncertainty correspondence in place of the fixed \(\mathcal P\). Such specifications must be identified as model assumptions, not inferred from the agenda. Likewise, state whether uncertain market-impact elasticities are varied across models or included in the uncertainty specification, and report the resulting scope of each robustness claim.

Specify a common numeraire, aggregation weights, horizon and discounting for \(C\) and every component of \(L\), together with an accounting map for the affected population. Distinguish execution shortfalls that are transfers from social resource losses; document exclusions preventing overlap among settlement, default and real-economy losses and between \(C\) and \(L\). These measurement choices are research deliverables; no particular weights or horizon are prescribed here. For empirical claims, declare observed variables, measurement and sampling coverage, the admissible structural model class and restrictions connecting observations to counterfactual architectural responses. Distinguish calibrated comparisons from identification, partial identification and estimation, and report sensitivity or identified sets where observations do not determine the target.

For the displayed objective, write \(T(a)=\sup_{P\in\mathcal P}E_P[L(a,\omega)]\). Report the attainable nondominated pairs \((C(a),T(a))\): a pair is dominated if another feasible architecture has no larger value in either coordinate and a strictly smaller value in at least one. Report limiting pairs separately when they are not attained. This is an expected-stress-loss frontier, not automatically a quantile or expected-shortfall frontier. Any additional tail-loss or resilience claim must declare its statistic or stress-event criterion, the relevant laws and horizon, and the evidence supporting it; a claim of meeting a threshold must state and justify that threshold. The agenda itself supplies no universal resilience threshold. An empirical optimization claim must explain how measurement, identification, estimation and optimization uncertainty affect the comparison or stated guarantee. Merely invoking the existence of approximate minimizers is not an empirical characterization of resilient designs.

\textbf{Additional formal definitions and scope (editorial).} The shock state \(\omega\) includes order-flow, collateral-value and funding innovations over a declared stress horizon. A settlement network specifies obligations, due dates, legally enforceable netting sets and the default waterfall; execution costs are measured relative to a stated pretrade benchmark. Loss \(L\) counts resource destruction and externalities without double-counting transfers between traders. Declare a common measurable shock space \((\Omega,\mathcal F)\). \(\mathcal P\) is a nonempty set of probability measures on this space, and for every \(a\in A\), \(L(a,\cdot)\) is measurable with \(E_P[|L(a,\omega)|]<\infty\) for every \(P\in\mathcal P\). These individual integrability conditions do not imply \(T(a)<\infty\); allow \(J(a)=+\infty\) and report architectures with infinite robust loss separately from an empty feasible set. For nonempty compact \(A\) in a declared topology, continuous finite \(C\), and continuous finite \(a\mapsto E_P[L(a,\omega)]\) for each \(P\), if \(J\) is finite on \(A\), it is lower semicontinuous and a minimizer exists. Robust finiteness is an additional hypothesis, not a consequence of individual integrability. That existence is a standard result; the research questions are identification of costs and tail losses, economically justified restrictions, and comparative characterization of the minimizing designs. For nonempty \(A\) and finite \(J_*\), the design target is, for each specified \(\varepsilon>0\), a feasible architecture \(a_\varepsilon\in A\) satisfying \(J(a_\varepsilon)\leq J_*+\varepsilon\), with the supporting comparison justified under the declared model and uncertainty specification. Report infeasibility or an unbounded value separately. An exact optimizer is requested only under additional attainment assumptions; it need not exist for the general rule class.
\subsection*{Variants and assumption changes}
\begin{enumerate}
\item \textbf{Known stress law (editorial).} Fix \(P\) and institution behavior primitives, then compare bilateral settlement, netting and central clearing using a common loss.
\item \textbf{Ambiguous stress dependence (editorial).} Vary joint funding/selling shocks over \(\mathcal P\), allowing leverage responses to rules; identify the routine-cost/tail-loss frontier.
\end{enumerate}
\subsection*{References and source audit}
\begin{enumerate}
\item Official January 9, 2026 web version; locator: The Financial System / Financial market dynamics; 2026 core-market priority. Broad agenda; exact model editorial. URL: \url{https://www.bankofengland.co.uk/research/bank-of-england-agenda-for-research}.
\item BIS WP 972 cover verifies authors; October 2021 original, January 27, 2022 revision inspected; no DOI assigned. Locator: Section 5, printed pp.38--39 (PDF pp.41--42). Support: Nonbank regulation/backstop agenda. Full January 2022 revision accessible. URL: \url{https://www.bis.org/publications/working-paper-972-non-bank-financial-intermediaries-and-financial-stability.pdf}.
\end{enumerate}
\begin{enumerate}\raggedright
\item\hypertarget{definitions:source:30007599:1}{} Bank of England. 2026. \emph{Bank of England Agenda for Research, 2025–2028}. The Financial System / Financial market dynamics; 2026 core-market priority.
\url{https://www.bankofengland.co.uk/research/bank-of-england-agenda-for-research}.
\item\hypertarget{definitions:source:30007599:2}{} Sirio Aramonte; Andreas Schrimpf; Hyun Song Shin. 2021. \emph{Non-bank financial intermediaries and financial stability}. Section 5, printed pp.38--39 (PDF pp.41--42).
\url{https://www.bis.org/publications/working-paper-972-non-bank-financial-intermediaries-and-financial-stability.pdf}.
\end{enumerate}

\subsection*{Common conventions: Source mathematical conventions}

These conventions apply to each entry unless explicitly replaced.

\textbf{Models and identification.} A primitive model \(m\) specifies all probability laws, feasible choices, information, equilibrium and measurement rules used by its target. The admissible class \(\mathcal M\) is fixed independently of outcomes. If \(P_O(m)\) is its observable probability law and \(\tau(m)\) the target, its identified set at observable law \(P\) is
\[
 \mathcal I_\tau(P)=\{\tau(m):m\in\mathcal M,\ P_O(m)=P\}.
\]
Point identification means that this set is a singleton. An empty set signals incompatibility of data and assumptions. Coordinate bounds need not describe the joint set. Bounds are sharp only if every retained target is attainable or arbitrarily approachable by compatible models. Finite bounds require explicit restrictions on unobserved quantities; unrestricted confounding, hidden wealth, extrapolation or arbitrary equilibrium selection can yield uninformative sets.

\textbf{Causal interventions.} A potential outcome \(Y_i(p)\) is the outcome under a complete policy regime \(p\), including timing, financing, information and market spillovers. The target population and units are fixed before treatment. With interference, use the complete assignment vector or a justified exposure mapping; individual no-interference notation alone cannot identify an economy-wide intervention. A difference of mean potential outcomes does not require the cross-world joint distribution. Conditioning on post-treatment migration, survival, enrollment or employment changes the estimand and needs separate justification.

\textbf{Statistical statements.} An estimator, confidence set, forecast or test specifies a sampling law, model class and loss/coverage criterion. Uniform \((1-\alpha)\)-coverage means \(\inf_{m\in\mathcal M}\Pr_m\{\tau(m)\in C_n\}\geq1-\alpha\), or an explicitly asymptotic version. The whole parameter space trivially has coverage, so informativeness requires a stated width, risk or power objective. A forecasting improvement is relative to a named benchmark, information set and evaluation population.

\textbf{Equilibrium and optimization.} An equilibrium comprises feasible plans, optimality, beliefs where needed, budget constraints and all market/resource clearing. The primitives or a stated benchmark define each correspondence \(\mathcal E(p;m)\). Nonemptiness is assumed or proved, never inferred just from compact policy sets. A selected-equilibrium target uses a declared selection rule; a robust target quantifies over all permitted equilibria. Use suprema or infima unless attainment is established. An \(\varepsilon\)-optimal policy is feasible and within \(\varepsilon\) of the finite optimum value. Report an empty feasible set separately.

\textbf{Welfare and accounting.} Normative utility, interpersonal normalization, social weights, numeraire, discounting and the use of public receipts are primitives. Behavioral utility need not coincide with the welfare criterion. Household welfare includes ownership income and rents once; adding profits again to compensating variation generally double-counts them. A monetary equivalent is defined by an explicit transfer and price convention, with monotonicity and range conditions. Average effects, mechanism contributions, transfers and welfare losses are different targets. Infinite sums require convergence and dynamic feasible plans require terminal or no-Ponzi conditions.

\textbf{Source versions.} A working-paper date, revision date and journal publication year are distinguished. A DOI for a journal version is not automatically the DOI of an earlier manuscript. Locators refer to the stated version and printed pages where specified. A metadata-only check does not verify a claim about a full-text conclusion. Every entry's source subsection narrows the provenance claim accordingly.
% END ECONOMICS STATEMENT
\end{document}
